\section{Simple linear regression}

The entries in this section use one invented data set with a declared sampling model. Ten corporate bonds have leverage $x=20, 25, \dots, 65$ (debt over assets, in percent), fixed by design. Their credit spreads $y$ (basis points) were generated as $y=40+5x+\varepsilon$, with independent normal errors of standard deviation 30, and rounded to whole basis points: $102, 216, 229, 240, 222, 197, 265, 326, 388, 349$. The generating line is known here so that estimates can be compared with it; in practice it is unknown. The data illustrate the mechanics and are not evidence about credit markets.

\begin{qfmethod}[Simple linear regression by least squares: the population line, the estimated intercept and slope, fitted values and residuals, coefficient units, and association against causation]\label{meth:simple-regression-estimation}
\qffield{Statement} Simple linear regression relates a response $y$ to one predictor $x$. The \emph{population regression line} $\mu_Y=E(Y)=\beta_0+\beta_1x$ connects the mean responses of the subpopulations at each $x$; a sample gives the estimated line $\hat y_i=b_0+b_1x_i$, and another sample would give a different one \citep[sec.~1.3]{regression-methods-course-notes}. The prediction error, or residual, is $e_i=y_i-\hat y_i$. The \emph{least squares criterion} chooses $b_0$ and $b_1$ to ``minimize the sum of the squared prediction errors'' $Q=\sum_{i=1}^n(y_i-\hat y_i)^2$, which gives
\[ b_1=\frac{\sum_{i=1}^n(x_i-\bar x)(y_i-\bar y)}{\sum_{i=1}^n(x_i-\bar x)^2},\qquad b_0=\bar y-b_1\bar x \]
(sec.~1.2). ``Association is not causation'': a large $R^2$ ``does not imply that $x$ causes the changes in $y$'' (sec.~1.5). Experiments, in which the researcher controls the predictor, support causal conclusions; observational studies, the source of most regression data, do not (sec.~1.8, Caution~4).

\qffield{Notation} $y$ is the response (dependent variable), $x$ the predictor (independent or explanatory variable), $n$ the number of observations. $b_1$ has the units of $y$ per unit of $x$, here basis points of spread per percentage point of leverage; $b_0$ is the fitted value at $x=0$, in units of $y$. Interpretation: the first-order conditions of the criterion make the residuals sum to zero and be uncorrelated with $x$ in the sample, and put the point $(\bar x,\bar y)$ on the fitted line.

\qffield{Conditions} The intercept and slope exist and are unique when the data are finite and $S_{xx}=\sum(x_i-\bar x)^2>0$, that is, when at least two predictor values differ. If every $x_i$ is the same, every line through $(\bar x,\bar y)$ minimises $Q$ and the slope is not identified. This is the full-rank condition on the design matrix \citep[Normal linear regression model, ``Setting'' and ``Assumptions'']{linear-regression-lectures}, written for one predictor (Interpretation). Fitting the line needs no assumption about the errors, normality included; the assumptions behind unbiasedness, standard errors and tests are set out in \hyperref[meth:regression-diagnostics-anova]{regression diagnostics}.

Reading $b_1$ as the average difference in $y$ between observations one unit of $x$ apart assumes the mean of $y$ is linear in $x$ over the observed range. Reading it as the effect of \emph{changing} $x$ needs a research design that rules out other explanations, which a cross-section of bonds does not supply: issuers with high leverage may differ in sector, size or asset risk (Interpretation). One unusual observation can move the slope substantially, so the data should be plotted \citep[sec.~1.8, Caution~3]{regression-methods-course-notes}. Regressing $x$ on $y$ gives a different line, not the inverse of this one (sec.~2.8, where the two directions give slopes 0.967 and 0.911).

\qffield{Practice} Example, the ten bonds: $\bar x=42.5$, $\bar y=253.4$, $\sum(x_i-\bar x)^2=2{,}062.5$ and $\sum(x_i-\bar x)(y_i-\bar y)=9{,}905$, so $b_1=9{,}905/2{,}062.5=4.80$ basis points per percentage point and $b_0=253.4-4.80\times42.5=49.3$. The bond at 45 percent leverage has a fitted spread of 265.4 and a residual of $197-265.4=-68.4$. In this sample, bonds with one percentage point more leverage have spreads about 4.8 basis points wider on average, close to the generating slope of 5. The intercept, a spread at zero leverage, lies far outside the data and has no practical meaning here.

\qffield{Pitfalls} Reading the slope as the effect of an issuer raising its leverage. Interpreting the intercept at an $x$ far from the data. Inverting a regression of $y$ on $x$ to predict $x$ from $y$. Fitting without plotting the data.

\qffield{See also} \hyperref[meth:regression-diagnostics-anova]{regression diagnostics}; \hyperref[meth:regression-prediction-functional-form]{regression predictions}; \hyperref[def:return-dependence]{covariance and correlation}; applications: \hyperref[mod:minimum-variance-hedge-ratio]{minimum-variance hedge ratio} (an OLS slope), \hyperref[meth:spread-volatility-regression]{spread--volatility regression}.

\end{qfmethod}

\begin{qfmethod}[Regression assumptions and diagnostics: what estimation, standard errors and exact tests each assume, residual plots, the standard error of estimate, $R^2$, coefficient $t$ tests and the simple-regression ANOVA, with $F=t^2$ for a zero slope]\label{meth:regression-diagnostics-anova}
\qffield{Statement} The simple linear regression model has four conditions: the mean of $Y_i$ is a \emph{linear} function of $x_i$; the errors are \emph{independent}; the errors are \emph{normally} distributed at each $x_i$; and they have \emph{equal} variance $\sigma^2$ \citep[sec.~1.3]{regression-methods-course-notes}. The mean square error $MSE=\sum(y_i-\hat y_i)^2/(n-2)$ estimates $\sigma^2$, and $S=\sqrt{MSE}$ (sec.~1.4); it needs $n>2$, because with two observations the line passes through both, $SSE=0$ and no degrees of freedom remain for $\sigma^2$, the standard errors or the tests (Interpretation). With $SSR=\sum(\hat y_i-\bar y)^2$, $SSE=\sum(y_i-\hat y_i)^2$ and $SSTO=SSR+SSE=\sum(y_i-\bar y)^2$, $R^2=SSR/SSTO=1-SSE/SSTO$ (sec.~1.5). The slope test uses $t^*=(b_1-\beta)/se(b_1)$, $se(b_1)=\sqrt{MSE}/\sqrt{\sum(x_i-\bar x)^2}$, on $n-2$ degrees of freedom, and the interval $b_1\pm t_{(\alpha/2,n-2)}\,se(b_1)$ (sec.~2.1).

The ANOVA table has 1, $n-2$ and $n-1$ degrees of freedom for regression, error and total; $MSR=SSR/1$, and $F^*=MSR/MSE$ is compared with an $F$ distribution with 1 and $n-2$ degrees of freedom (sec.~2.6). This $F^*$ tests $H_0:\beta_1=0$. For one predictor, with $t^*$ computed for that same null, $(t^*_{(n-2)})^2=F^*_{(1,n-2)}$: the $F$ test is ``algebraically equivalent to the two-tailed t-test'' and gives the same $p$-value, but only the $t$ test can test a one-sided alternative (sec.~2.7). A $t$ statistic for another value, $H_0:\beta_1=\beta$, squares to the $F$ statistic of that restriction in the general test of linear restrictions, not to the ANOVA $F^*$ \citep[Linear regression -- hypothesis testing, ``Test of a restriction on a single coefficient (t test)'' and ``Test of a set of linear restrictions (F test)'']{linear-regression-lectures}; the equality for one restriction is algebra (Interpretation).

Residuals estimate the unknown errors. A residuals-versus-fits plot detects non-linearity, unequal error variances and outliers; a systematic pattern suggests a non-linear regression function, and fanning suggests unequal variances (secs.~4.2 and~4.4). Standardised residuals beyond $\pm2$ are ``flagged for further investigation'', not removed (sec.~4.4). A residuals-versus-order plot detects serial correlation and ``is only appropriate if you know the order in which the data were collected'' (sec.~4.5); a normal probability plot checks normality (sec.~4.6). Sensitivity differs by condition: tests and intervals are ``very sensitive to even minor departures from independence'', sensitive to moderate departures from equal variance, fairly robust to non-normality for the coefficients, while prediction intervals are ``quite sensitive to departures from normality'' (sec.~4.1).

\qffield{Notation} $n-2$ degrees of freedom because two parameters are estimated. $S$ is the \emph{standard error of the estimate} (of the regression), in units of $y$. Interpretation: $se(b_0)=S\sqrt{1/n+\bar x^2/\sum(x_i-\bar x)^2}$; $F^*=R^2(n-2)/(1-R^2)$; with one predictor $R^2$ equals the squared sample correlation of $x$ and $y$.

\qffield{Conditions} Each conclusion needs a different subset of the conditions. Computing $b_0$ and $b_1$ needs only finite data with $S_{xx}>0$ (\hyperref[meth:simple-regression-estimation]{least squares estimation}). They are unbiased when, in addition, the mean of $Y$ is linear in $x$, so that $E(\varepsilon\mid X)=0$. Uncorrelated errors of equal variance, $\mathrm{Var}(\varepsilon\mid X)=\sigma^2I$, add the conventional standard errors above and make least squares the best linear unbiased estimator \citep[Gauss--Markov theorem, ``Assumptions'', ``OLS is linear and unbiased'' and ``The covariance matrix of the OLS estimator'']{linear-regression-lectures}. Normal errors add the exact $t$ and $F$ distributions on $n-2$ degrees of freedom (Normal linear regression model, ``Distribution of the OLS estimator'' and ``Estimation of the variance of the error terms''; Linear regression -- hypothesis testing, ``Tests of hypothesis in the normal linear regression model'').

Without normality the same statistics are large-sample approximations, justified when the estimator is asymptotically normal (Linear regression -- hypothesis testing, ``Tests of hypothesis when the OLS estimator is asymptotically normal''); ten observations give little assurance that the approximation is close. With independent, identically distributed errors of mean zero and finite variance, least squares is consistent for the coefficients whether or not the errors are normal, ``if the design matrix is asymptotically nondegenerate''; normality then affects only efficiency \citep[sec.~2, pp.~605--606]{smearing-retransformation-1983}. Full rank in every sample is not enough. With $x_i=2^{-i}$, $S_{xx}$ tends to $1/3$ as $n$ grows, so $\mathrm{Var}(b_1)=\sigma^2/S_{xx}$ tends to $3\sigma^2$, not zero; with one predictor, consistency of $b_1$ needs $S_{xx}\to\infty$ (Interpretation). Correlated errors invalidate the conventional standard errors and tests even when the coefficients remain unbiased (Interpretation).

A residual plot can reveal a violation; it cannot prove that the conditions hold. With ten points a curve, a fan or serial correlation can easily go unseen, and over-reading small plots is a known error \citep[sec.~4.2]{regression-methods-course-notes}. A large $R^2$ ``should not be interpreted as meaning that the estimated regression line fits the data well'', a significant slope need not be meaningfully different from zero, and a large $R^2$ does not guarantee useful intervals (sec.~1.8, Cautions~2, 6 and~7). $R^2$ and the tests describe the estimation sample; out-of-sample predictive performance needs separate evidence (Interpretation; \hyperref[pit:parameter-overfitting]{parameter overfitting}). Bonds of one issuer or sector, or observations in time order, may have dependent errors; time-series errors need time-series methods (sec.~4.5).

\qffield{Practice} Example, the ten bonds: $SSE=14{,}356.4$, $SSR=47{,}568.0$ and $SSTO=61{,}924.4$; $S=\sqrt{14{,}356.4/8}=42.4$ basis points and $R^2=0.768$. The slope has $se(b_1)=42.36/\sqrt{2{,}062.5}=0.933$ and $t=4.802/0.933=5.15$ on 8 degrees of freedom, $p=0.0009$; its 95 percent interval is 2.65 to 6.95. Testing $\beta_1=5$, the generating value, gives $t=-0.212$, not rejected; its square, 0.045, is the $F$ statistic for the restriction $\beta_1=5$, not the ANOVA's 26.51. The intercept has $t=1.18$, $p=0.27$.

\begin{center}\small
\begin{tabular}{@{}lrrrr@{}}
\toprule
Source & DF & SS & MS & $F$ \\
\midrule
Regression & 1 & 47{,}568.0 & 47{,}568.0 & 26.51 \\
Error & 8 & 14{,}356.4 & 1{,}794.5 & \\
Total & 9 & 61{,}924.4 & & \\
\bottomrule
\end{tabular}
\end{center}

$F=47{,}568.0/1{,}794.5=26.51=5.149^2$, the square of the $t$ statistic for a zero slope, with the same $p$-value as that two-sided $t$ test. The largest standardised residual is $-1.61$, at 45 percent leverage; none exceeds 2. These checks find no violation, which with ten bonds is weak evidence that none exists. Counterexample, original: the exact curve $y=0.2(x-10)^2$ on the same leverages gives a straight-line $R^2$ of 0.96, yet its residuals are positive, then negative, then positive.

\qffield{Pitfalls} Treating a clean residual plot as proof of the assumptions. Reading a high $R^2$ as a good model or as predictive power. Deleting observations because their residuals are large. Using the $F$ test for a one-sided question. Equating the ANOVA $F$ with the square of a $t$ statistic for a non-zero null value. Applying the $t$ test to serially correlated errors.

\qffield{See also} \hyperref[meth:simple-regression-estimation]{least squares estimation}; \hyperref[meth:hypothesis-testing]{hypothesis tests}; \hyperref[meth:regression-prediction-functional-form]{regression predictions}; \hyperref[meth:purged-embargoed-cross-validation]{purged cross-validation} (out-of-sample evaluation).

\end{qfmethod}

\begin{qfmethod}[Regression predictions: fitted values, intervals for the mean response and for a new observation, extrapolation, and log functional forms with retransformation to geometric means, medians and means]\label{meth:regression-prediction-functional-form}
\qffield{Statement} At a predictor value $x_h$, the confidence interval for the mean response is
\[ \hat y_h\pm t_{(1-\alpha/2,n-2)}\sqrt{MSE\Big(\frac1n+\frac{(x_h-\bar x)^2}{\sum(x_i-\bar x)^2}\Big)}, \]
narrower when $x_h$ is near $\bar x$; it applies when $x_h$ lies ``within the range of the x values in the data set'' and the LINE conditions hold, and tolerates approximately normal errors \citep[sec.~3.2]{regression-methods-course-notes}. The prediction interval for a new response adds an extra $MSE$ under the root, $MSE\,(1+1/n+(x_h-\bar x)^2/\sum(x_i-\bar x)^2)$: the variance of a prediction is the variance $\sigma^2$ of a new response about its mean plus the variance of $\hat y_h$. The prediction interval is therefore always wider and ``depends strongly on the condition that the error terms are normally distributed'' (sec.~3.3).

Transforming $x$ suits non-linearity alone; transforming $y$ suits non-normality or unequal variances (secs.~9.1--9.2). Predictions from a log response estimate $E(\ln Y\mid x)$ and come back in log units. In a log--level model a unit of $x$ multiplies the back-transformed value by $e^{\beta_1}$; in a log--log model a $k$-fold change in $x$ multiplies it by $k^{\beta_1}$. The course notes read the back-transformed value as the median, because the mean of the logs ``approximately equals the natural log of the median'' (secs.~9.2--9.3).

Exactly, $\exp E(\ln Y\mid x)$ is the conditional \emph{geometric mean} of $Y$. It is the conditional median only when the log-scale error has median zero, as a symmetric error does; Hyndman reads a back-transformed forecast as the median ``assuming the forecast densities on the transformed scale are symmetric'' \citep[para.~1]{back-transformed-forecast-mean-2014}. A skewed error with mean zero separates the two (Interpretation). When the error distribution does not change with $x$, the median and the mean change by the same factors as the geometric mean.

The mean is not the retransformed prediction: $E\,Y_0=E\,h(x_0\beta+\varepsilon)\ne h(x_0\beta)$. With normal errors and a log transformation, $E\,Y_0=\exp(x_0\beta+\sigma^2/2)$, estimated by $\exp(x_0\hat\beta+\hat\sigma^2/2)$. Under the design and moment conditions of \hyperref[meth:regression-diagnostics-anova]{regression diagnostics}, whatever the error distribution, that estimate is consistent for $\exp(x_0\beta+\sigma^2/2)$, which equals $E\,Y_0$ under normal errors but otherwise ``might not'': in Duan's normal-mixture example it has a 41 percent asymptotic bias. Normality is sufficient for it to estimate $E\,Y_0$ consistently, not necessary; what matters is whether $E\,e^{\varepsilon}=e^{\sigma^2/2}$ (Interpretation). The \emph{smearing estimate} replaces the error distribution by the residuals: $\frac1n\sum_{i=1}^nh(x_0\hat\beta+\hat\varepsilon_i)$ \citep[secs.~2--3, eq.~(3.1), pp.~605--606]{smearing-retransformation-1983}. It is weakly consistent for $E\,Y_0$ when $h$ is continuously differentiable, the design contains an intercept, $X'X/n$ converges to a positive-definite matrix and a moment condition holds, which for the log transformation is $E\,e^{2\varepsilon}<\infty$ (Theorem~1 and the text after it, sec.~4, p.~607). For a lognormal variable the median is the scale parameter and the mean exceeds it by the factor $e^{0.5\sigma^2}$ \citep[sec.~1.3.6.6.9]{statistical-methods-handbook}.

\qffield{Notation} Interpretation: the \emph{standard error of the estimate} $S=\sqrt{MSE}$ measures scatter about the line; the standard error of the fit is $S\sqrt{1/n+(x_h-\bar x)^2/\sum(x_i-\bar x)^2}$; the standard error of prediction is $\sqrt{S^2+(\text{standard error of the fit})^2}$. Forms and slope units:
\begin{itemize}
\item level--level, $y=b_0+b_1x$: $b_1$ units of $y$ per unit of $x$;
\item log--level, $\ln y=b_0+b_1x$: the geometric mean of $y$ changes by $100(e^{b_1}-1)$ percent per unit of $x$;
\item level--log, $y=b_0+b_1\ln x$: $y$ changes by $b_1\ln(1.01)\approx b_1/100$ units when $x$ rises 1 percent;
\item log--log, $\ln y=b_0+b_1\ln x$: $b_1$ is an elasticity; a 1 percent rise in $x$ changes the geometric mean of $y$ by $100(1.01^{b_1}-1)\approx b_1$ percent, a fraction $1.01^{b_1}-1$.
\end{itemize}
Logarithms need positive values: spreads, prices and volumes qualify, returns do not.

\qffield{Conditions} The intervals assume the fitted model is correct where they are used. Outside the observed $x$ the formula widens only for parameter uncertainty; it says nothing about whether the line still holds, which no data in the sample can check (Interpretation). A prediction interval is for one new observation from the same population, with an error independent of the sample. $R^2$ values of models with $y$ and with $\ln y$ measure different responses and cannot rank the forms. The normal-theory mean correction needs normal, equal-variance errors on the log scale; the smearing estimate drops normality but still treats the errors as independent draws from one distribution, so neither corrects for variance that changes with $x$ (Interpretation, from Duan's model assumptions).

\qffield{Practice} Example, the ten bonds at 45 percent leverage: $\hat y=49.3+4.802\times45=265.4$ basis points. The standard error of the fit is 13.60 and of prediction $\sqrt{42.36^2+13.60^2}=44.49$. With $t_{(0.975,8)}=2.306$, the 95 percent interval for the mean spread of all such bonds is 234.1 to 296.8. For one new bond it is 162.8 to 368.0. In 20{,}000 samples simulated from the declared model, the first interval covered the true mean of 265 in 94.8 percent of samples and the second covered a new bond's spread in 95.0 percent; the mean interval misused for a new bond covered it only 49.7 percent of the time (49.9 percent exactly).

Extrapolating to 90 percent leverage gives 481.5 with a standard error of the fit of 46.3, 3.4 times that at 45 percent, and a prediction interval of 336.8 to 626.2. That interval is valid only if the straight line, its variance and its normal errors extend to leverage no bond in the sample has.

The log--level fit to the same bonds is $\ln y=4.597+0.02075x$ with $S^2=0.0472$: each percentage point of leverage multiplies the geometric-mean spread by $e^{0.02075}=1.021$, 2.1 percent. At 45 percent the fitted geometric mean is $e^{5.530}=252.2$, also the median if the log errors have median zero; the normal-theory mean is $252.2\times e^{0.0472/2}=258.2$ and the smearing estimate is 256.8. The log--log elasticity is 0.826: a 1 percent rise in leverage raises the geometric-mean spread by $100(1.01^{0.826}-1)=0.825$ percent, a fraction 0.00825. The level--log slope of 183.9 adds 1.83 basis points per 1 percent rise in leverage. The data came from the level model, so these fits show the arithmetic of each form; their $R^2$ values, 0.70 to 0.75, are not evidence against the level model's 0.77.

Counterexample, original: let the log error be $\varepsilon=0.5(E-1)$ with $E$ standard exponential. Its mean is zero and its variance 0.25, but it is right-skewed, with median $0.5(\ln2-1)=-0.153$. Relative to $e^{x\beta}$, the geometric mean of $Y$ is 1, the median 0.858, the mean $e^{-0.5}/0.5=1.213$ and the normal-theory factor $e^{0.25/2}=1.133$. Exponentiating the log prediction then overstates the median by 16.6 percent and understates the mean by 17.6 percent; the normal-theory correction is still 6.6 percent below the mean.

\qffield{Pitfalls} Quoting a confidence interval for the mean as the range for one new outcome. Extrapolating beyond the observed predictors. Exponentiating a log prediction and reporting it as the expected value, or as the median when the log errors may be skewed. Reading the fraction $1.01^{b_1}-1$ as a percentage. Choosing a functional form by comparing $R^2$ across different responses. Taking logs of returns or other quantities that can be zero or negative.

\qffield{See also} \hyperref[meth:regression-diagnostics-anova]{regression diagnostics}; \hyperref[meth:confidence-intervals]{confidence intervals}; \hyperref[def:financial-distributions]{lognormal moments}; \hyperref[meth:spread-volatility-regression]{spread--volatility regression} (a log--log application with the lognormal correction).

\end{qfmethod}
